\section{Local Invariants}

In commuting coordinates $(x^1,x^2)$, suppose a vector field $Z$ is obtained
by parallel transport first along the $x^1$-axis and then along the
$x^2$-coordinate lines. It satisfies $\nabla_{\partial_2}Z=0$, and it is
parallel precisely when one can propagate
\[
\nabla_{\partial_2}\nabla_{\partial_1}Z=0. \qquad (7.1)
\]
The relevant commutation identity is
\begin{equation}
\nabla_{\partial_2}\nabla_{\partial_1}Z
=\nabla_{\partial_1}\nabla_{\partial_2}Z. \tag{7.2}
\end{equation}
For the Euclidean connection and arbitrary vector fields,
\[
\overline\nabla_X\overline\nabla_YZ
-\overline\nabla_Y\overline\nabla_XZ
=\overline\nabla_{[X,Y]}Z,
\]
because $XY(Z^k)-YX(Z^k)=[X,Y](Z^k)$.

A Riemannian manifold is \emph{flat} if every point has a neighborhood
isometric to an open subset of Euclidean space. A pseudo-Riemannian manifold is
flat when the corresponding local model is pseudo-Euclidean space. A
connection $\nabla$ satisfies the \emph{flatness criterion} if, for all vector
fields on every open subset,
\begin{equation}
\nabla_X\nabla_YZ-\nabla_Y\nabla_XZ=\nabla_{[X,Y]}Z. \tag{7.3}
\end{equation}

\begin{example}
\label{ex:n-torus-flat-metric}
\uses{ex:n-torus-riemannian-submanifold}
\dcref{ch7:7.1}

The metric on the $n$-torus induced by the embedding in $\mathbb R^{2n}$ from
Example 2.21 is flat: every point has coordinates in which the metric is
Euclidean.
\end{example}

\begin{proposition}
\label{prop:flat-manifold-levi-civita-flatness}
\dcref{ch7:7.2}
The Levi-Civita connection of every flat Riemannian or pseudo-Riemannian
manifold satisfies the flatness criterion.
\end{proposition}

\begin{proof}
\sketch The Euclidean and pseudo-Euclidean connections satisfy (7.3) by the
coordinate calculation above. Naturality of the Levi-Civita connection
transfers this identity through each local isometry to a model space.
\end{proof}

\section{The Curvature Tensor}

Let $(M,g)$ be Riemannian or pseudo-Riemannian. With the sign convention used
here, define
\[
R(X,Y)Z
=\nabla_X\nabla_YZ-\nabla_Y\nabla_XZ-\nabla_{[X,Y]}Z.
\]

\begin{proposition}
\label{prop:curvature-multilinearity}
\uses{lem:tensor-characterization-lemma}
\dcref{ch7:7.3}

The map $R$ is multilinear over $C^\infty(M)$ and therefore defines a
$(1,3)$-tensor field.
\end{proposition}

\begin{proof}
\sketch Real multilinearity is immediate. For $f\in C^\infty(M)$,
\begin{align*}
R(X,fY)Z
&=\nabla_X(f\nabla_YZ)-f\nabla_Y\nabla_XZ
  -\nabla_{f[X,Y]+(Xf)Y}Z\\
&=fR(X,Y)Z.
\end{align*}
Skew-symmetry in $X,Y$ gives $C^\infty(M)$-linearity in $X$. Expanding
$R(X,Y)(fZ)$, the first-derivative terms cancel, and the scalar second-order
term is
$XYf-YXf-[X,Y]f=0$; hence $R$ is also $C^\infty(M)$-linear in $Z$.
The tensor characterization lemma now applies.
\end{proof}

For vector fields $X,Y$, the bundle endomorphism $Z\mapsto R(X,Y)Z$ is the
\emph{curvature endomorphism}. The tensor $R$ is the $(1,3)$-curvature tensor.
Its contravariant index is placed last, contrary to the default ordering of
tensor arguments:
\[
R=R_{ijk}^{\ell}\,dx^i\otimes dx^j\otimes dx^k\otimes\partial_\ell,
\qquad
R(\partial_i,\partial_j)\partial_k=R_{ijk}^{\ell}\partial_\ell.
\]

\begin{proposition}
\label{prop:curvature-tensor-components}
\dcref{ch7:7.4}
For a Riemannian or pseudo-Riemannian manifold, the components of the
$(1,3)$-curvature tensor in any smooth coordinates are
\begin{equation}
R_{ijk}^l
=\partial_i\Gamma_{jk}^l-\partial_j\Gamma_{ik}^l
+\Gamma_{jk}^m\Gamma_{im}^l-\Gamma_{ik}^m\Gamma_{jm}^l. \tag{7.4}
\end{equation}
\end{proposition}

Let $\Gamma\colon J\times I\to M$ be a smooth family of curves. Write
$\partial_t\Gamma$ and $\partial_s\Gamma$ for its two velocity fields and
$D_t,D_s$ for the associated covariant derivatives along $\Gamma$.

\begin{proposition}
\label{prop:curvature-commutator-failure}
\uses{prop:connection-on-tensor-bundles}
\dcref{ch7:7.5}

For every smooth vector field $V$ along a smooth family
$\Gamma\colon J\times I\to M$ in a Riemannian or pseudo-Riemannian manifold,
\[
D_sD_tV-D_tD_sV=R(\partial_s\Gamma,\partial_t\Gamma)V.
\]
\end{proposition}

\begin{proof}
\sketch In coordinates, write
\[
\Gamma(s,t)=(y^1(s,t),\ldots,y^n(s,t)),
\qquad V=V^i\partial_i|_{\Gamma(s,t)}.
\]
Applying the coordinate formula for covariant differentiation twice and
subtracting cancels all derivatives of $V^i$ except
\begin{equation}
D_sD_tV-D_tD_sV
=V^i(D_sD_t\partial_i-D_tD_s\partial_i). \tag{7.6}
\end{equation}
Set
\[
S=\partial_s\Gamma=\frac{\partial y^k}{\partial s}\partial_k,
\qquad
T=\partial_t\Gamma=\frac{\partial y^j}{\partial t}\partial_j.
\]
Since the coordinate fields and their covariant derivatives are extendible,
\begin{align*}
D_sD_t\partial_i
&=\frac{\partial^2y^j}{\partial s\partial t}
  \nabla_{\partial_j}\partial_i
 +\frac{\partial y^j}{\partial t}
  \frac{\partial y^k}{\partial s}
  \nabla_{\partial_k}\nabla_{\partial_j}\partial_i.
\end{align*}
Interchanging $s,t$ and subtracting eliminates the mixed derivative term and
gives
\[
D_sD_t\partial_i-D_tD_s\partial_i
=R(S,T)\partial_i.
\]
Substitution in (7.6) proves the formula.
\end{proof}

The \emph{Riemann curvature tensor} is the covariant tensor
$\operatorname{Rm}=R^\flat$ obtained by lowering the last index of $R$:
\begin{equation}
\operatorname{Rm}(X,Y,Z,W)=\langle R(X,Y)Z,W\rangle_g. \tag{7.7}
\end{equation}
Thus in smooth coordinates
\[
\operatorname{Rm}=R_{ijkl}\,dx^i\otimes dx^j\otimes dx^k\otimes dx^l,
\qquad R_{ijkl}=g_{lm}R_{ijk}^m,
\]
and
\begin{equation}
R_{ijkl}=g_{lm}\left(
\partial_i\Gamma_{jk}^m-\partial_j\Gamma_{ik}^m
+\Gamma_{jk}^p\Gamma_{ip}^m-\Gamma_{ik}^p\Gamma_{jp}^m
\right). \tag{7.8}
\end{equation}
These formulas fix both conventions: $R$ has contravariant index last and the
sign is determined by the displayed commutator definition. Other conventions
may negate $R$ or permute the arguments used to define $\operatorname{Rm}$.

\begin{proposition}
\label{prop:curvature-isometry-invariant}
\dcref{ch7:7.6}
Let $(M,g)$ and $(\widetilde M,\widetilde g)$ be Riemannian or
pseudo-Riemannian manifolds. If
$\varphi\colon M\to\widetilde M$ is a local isometry, then
\[
\varphi^*\widetilde{\operatorname{Rm}}=\operatorname{Rm}.
\]
\end{proposition}
\begin{proof}
\sketch
The differential of a local isometry intertwines the two Levi--Civita
connections by uniqueness of the torsion-free metric connection. Applying
this identity twice shows that it intertwines the curvature endomorphisms;
pairing with the preserved metrics gives the displayed pullback identity.
\end{proof}

\section{Flat Manifolds}

\begin{lemma}
\label{lem:parallel-vector-field-flatness}
\uses{thm:linear-ode-existence-uniqueness}
\dcref{ch7:7.8}

Let $\nabla$ be a connection on a smooth manifold $M$ satisfying the flatness
criterion. For every $p\in M$ and $v\in T_pM$, there is a parallel vector field
$V$ on a neighborhood of $p$ such that $V_p=v$.
\end{lemma}

\begin{proof}
\sketch Choose coordinates centered at $p$ on a cube
\[
C_\varepsilon=\{x:|x^i|<\varepsilon,\ 1\leq i\leq n\}.
\]
Starting with $v$, successively parallel transport along coordinate lines in
the $x^1,x^2,\ldots,x^n$ directions. Smooth dependence for linear ODEs makes
the resulting vector field $V$ smooth. Let
\[
M_k=\{x\in C_\varepsilon:x^{k+1}=\cdots=x^n=0\}.
\]
Inductively prove
\begin{equation}
\nabla_{\partial_1}V=\cdots=\nabla_{\partial_k}V=0
\quad\text{on }M_k. \tag{7.9}
\end{equation}
The case $k=1$ follows from the construction. Assuming the assertion on
$M_k$, construction gives $\nabla_{\partial_{k+1}}V=0$ on $M_{k+1}$, and for
$i\leq k$ the flatness criterion and
$[\partial_{k+1},\partial_i]=0$ give
\[
\nabla_{\partial_{k+1}}(\nabla_{\partial_i}V)
=\nabla_{\partial_i}(\nabla_{\partial_{k+1}}V)=0.
\]
Thus $\nabla_{\partial_i}V$ is parallel along the $x^{k+1}$-lines and vanishes
where they meet $M_k$, so it vanishes on $M_{k+1}$. At $k=n$, $V$ is parallel
on the entire cube.
\end{proof}

\begin{theorem}
\label{thm:flatness-iff-zero-curvature}
\uses{prop:flat-manifold-levi-civita-flatness, lem:parallel-vector-field-flatness}
\dcref{ch7:7.10}

A Riemannian or pseudo-Riemannian manifold is flat if and only if its curvature
tensor vanishes identically.
\end{theorem}

\begin{proof}
\sketch If the metric is flat, Proposition 7.2 gives the flatness criterion,
so $R$ and $\operatorname{Rm}$ vanish.

Conversely, suppose $\operatorname{Rm}=0$. Nondegeneracy of $g$ implies
$R=0$, so the Levi-Civita connection satisfies the flatness criterion. At
$p\in M$, choose an orthonormal basis $(b_1,\ldots,b_n)$, in standard signature
order in the pseudo-Riemannian case. The preceding lemma extends it to parallel
fields $(E_1,\ldots,E_n)$ near $p$. Metric compatibility makes this a local
orthonormal frame. Since the Levi-Civita connection is torsion-free,
\[
[E_i,E_j]=\nabla_{E_i}E_j-\nabla_{E_j}E_i=0.
\]
The canonical form theorem for commuting vector fields provides local
coordinates $(y^1,\ldots,y^n)$ with $E_i=\partial/\partial y^i$. In these
coordinates,
\[
g_{ij}=\varepsilon_i\delta_{ij},
\qquad \varepsilon_i\in\{1,-1\},
\]
with the signs in standard order. Hence the coordinate map is a local isometry
to the Euclidean or appropriate pseudo-Euclidean model.
\end{proof}

\begin{theorem}
\label{thm:curvature-endomorphism-formula}
\uses{thm:linear-ode-existence-uniqueness, prop:curvature-commutator-failure}
\dcref{ch7:7.11}

Let $(M,g)$ be Riemannian or pseudo-Riemannian, let $I$ be an open interval
containing $0$, and let $\Gamma\colon I\times I\to M$ be smooth. Put
\[
p=\Gamma(0,0),
\qquad x=\partial_s\Gamma(0,0),
\qquad y=\partial_t\Gamma(0,0).
\]
For $s_1,s_2,t_1,t_2\in I$, let
\[
P_{s_1,t_1}^{s_1,t_2}
\colon T_{\Gamma(s_1,t_1)}M\longrightarrow T_{\Gamma(s_1,t_2)}M
\]
denote parallel transport along $t\mapsto\Gamma(s_1,t)$, and let
\[
P_{s_1,t_1}^{s_2,t_1}
\colon T_{\Gamma(s_1,t_1)}M\longrightarrow T_{\Gamma(s_2,t_1)}M
\]
denote parallel transport along $s\mapsto\Gamma(s,t_1)$. Then every
$z\in T_pM$ satisfies
\begin{equation}
R(x,y)z
=\lim_{\delta,\varepsilon\to0}
\frac{
P_{\delta,0}^{0,0}\circ
P_{\delta,\varepsilon}^{\delta,0}\circ
P_{0,\varepsilon}^{\delta,\varepsilon}\circ
P_{0,0}^{0,\varepsilon}(z)-z
}{\delta\varepsilon}. \tag{7.10}
\end{equation}
\end{theorem}

\begin{proof}
\sketch Define $Z$ along $\Gamma$ by parallel transporting $z$ first along
$t\mapsto\Gamma(0,t)$ and then, for each $t$, along
$s\mapsto\Gamma(s,t)$. Smooth dependence for linear ODEs makes $Z$ smooth,
and
\[
D_tZ(0,t)=0,
\qquad D_sZ(s,t)=0.
\]
Proposition~\ref{prop:curvature-commutator-failure} therefore gives
\[
R(x,y)z=D_sD_tZ(0,0)-D_tD_sZ(0,0)=D_sD_tZ(0,0).
\]
Covariant differentiation by parallel transport gives
\[
D_tZ(s,0)=\lim_{\varepsilon\to0}
\frac{P_{s,\varepsilon}^{s,0}(Z(s,\varepsilon))-Z(s,0)}{\varepsilon},
\qquad (7.11)
\]
and
\[
D_sD_tZ(0,0)=\lim_{\delta\to0}
\frac{P_{\delta,0}^{0,0}(D_tZ(\delta,0))-D_tZ(0,0)}{\delta}.
\qquad (7.12)
\]
Substitution yields
\[
D_sD_tZ(0,0)
=\lim_{\delta,\varepsilon\to0}
\frac{
P_{\delta,0}^{0,0}P_{\delta,\varepsilon}^{\delta,0}(Z(\delta,\varepsilon))
-P_{\delta,0}^{0,0}(Z(\delta,0))
-P_{0,\varepsilon}^{0,0}(Z(0,\varepsilon))
+Z(0,0)
}{\delta\varepsilon}.
\qquad (7.13)
\]
Parallelity in the two construction directions implies
\[
P_{\delta,0}^{0,0}(Z(\delta,0))
=P_{0,\varepsilon}^{0,0}(Z(0,\varepsilon))
=Z(0,0)=z
\]
and
\[
Z(\delta,\varepsilon)
=P_{0,\varepsilon}^{\delta,\varepsilon}(Z(0,\varepsilon))
=P_{0,\varepsilon}^{\delta,\varepsilon}
 \circ P_{0,0}^{0,\varepsilon}(z).
\]
Inserting these identities into (7.13) gives (7.10).
\end{proof}
\section{Symmetries of the Curvature Tensor}

\begin{proposition}[Symmetries of the Curvature Tensor]
\label{prop:curvature-tensor-symmetries}
\dcref{ch7:7.12}

Let $(M,g)$ be a Riemannian or pseudo-Riemannian manifold. For all vector
fields $W,X,Y,Z$, its $(0,4)$-curvature tensor satisfies
\begin{enumerate}[label=(\alph*)]
\item $Rm(W,X,Y,Z)=-Rm(X,W,Y,Z)$.
\item $Rm(W,X,Y,Z)=-Rm(W,X,Z,Y)$.
\item $Rm(W,X,Y,Z)=Rm(Y,Z,W,X)$.
\item $Rm(W,X,Y,Z)+Rm(X,Y,W,Z)+Rm(Y,W,X,Z)=0$.
\end{enumerate}
\end{proposition}

Identity (d) is the \emph{algebraic Bianchi identity}. In any local frame,
the symmetries are
\begin{align*}
R_{ijkl}&=-R_{jikl},&
R_{ijkl}&=-R_{ijlk},\\
R_{ijkl}&=R_{klij},&
R_{ijkl}+R_{jkil}+R_{kijl}&=0.
\end{align*}

\begin{proof}
\sketch Part (a) follows from $R(W,X)=-R(X,W)$. For (b), metric
compatibility gives
\begin{align*}
WX|Y|^2&=2\langle\nabla_W\nabla_XY,Y\rangle
 +2\langle\nabla_XY,\nabla_WY\rangle,\\
XW|Y|^2&=2\langle\nabla_X\nabla_WY,Y\rangle
 +2\langle\nabla_WY,\nabla_XY\rangle,\\
[W,X]|Y|^2&=2\langle\nabla_{[W,X]}Y,Y\rangle.
\end{align*}
Subtracting the last two equations from the first yields
$Rm(W,X,Y,Y)=0$; polarization in $Y,Z$ proves (b).

For (d), symmetry of the connection expands the cyclic sum as
\begin{align*}
R(W,X)Y+R(X,Y)W+R(Y,W)X
&=[W,[X,Y]]+[X,[Y,W]]+[Y,[W,X]],
\end{align*}
which vanishes by the Jacobi identity. Finally, apply (d) to the four cyclic
arrangements of $W,X,Y,Z$, add the resulting equations, and use (a) and (b);
the uncancelled terms give
$2Rm(Y,W,X,Z)-2Rm(X,Z,Y,W)=0$, which is equivalent to (c).
\end{proof}

\begin{proposition}[Differential Bianchi Identity]
\label{prop:differential-bianchi-identity}
\uses{prop:properties-normal-coordinates}
\dcref{ch7:7.13}


On a Riemannian or pseudo-Riemannian manifold,
\begin{equation}
\nabla Rm(X,Y,Z,V,W)+\nabla Rm(X,Y,V,W,Z)
+\nabla Rm(X,Y,W,Z,V)=0.
\tag{7.14}
\end{equation}
Equivalently, in components,
\begin{equation}
R_{ijkl;m}+R_{ijlm;k}+R_{ijmk;l}=0.
\tag{7.15}
\end{equation}
\end{proposition}

\begin{proof}
\sketch Curvature symmetries reduce (7.14) to
\begin{equation}
\nabla Rm(Z,V,X,Y,W)+\nabla Rm(V,W,X,Y,Z)
+\nabla Rm(W,Z,X,Y,V)=0.
\tag{7.16}
\end{equation}
Fix $p$ and use normal coordinate vector fields $X,Y,Z,V,W$ at $p$. Their
brackets vanish identically and their covariant derivatives vanish at $p$.
Thus
\[
(\nabla_WRm)(Z,V,X,Y)
=\langle\nabla_W\nabla_Z\nabla_VX-\nabla_W\nabla_V\nabla_ZX,Y\rangle
\quad\text{at }p.
\]
Adding its cyclic permutations in $W,Z,V$ gives
\[
\left\langle
R(W,Z)(\nabla_VX)+R(Z,V)(\nabla_WX)+R(V,W)(\nabla_ZX),Y
\right\rangle=0
\]
at $p$. Multilinearity and the arbitrariness of $p$ prove the identity.
\end{proof}

\section{The Ricci Identities}

For a tensor field $F$ of type $(k,l)$, write
$\nabla^2F=\nabla(\nabla F)$ and
$\nabla^2_{X,Y}F=\nabla^2F(\ldots,Y,X)$. Define the dual curvature map by
\[
(R(X,Y)^*\eta)(Z)=\eta(R(X,Y)Z).
\]

\begin{theorem}[Ricci Identities]
\label{thm:ricci-identities}
\dcref{ch7:7.14}

On a Riemannian or pseudo-Riemannian manifold, the following identities hold.
For a vector field $Z$,
\[
\nabla^2_{X,Y}Z-\nabla^2_{Y,X}Z=R(X,Y)Z.
\qquad (7.17)
\]
For a $1$-form $\beta$,
\[
\nabla^2_{X,Y}\beta-\nabla^2_{Y,X}\beta=-R(X,Y)^*\beta.
\qquad (7.18)
\]
For a $(k,l)$-tensor field $B$,
\begin{align}
&(\nabla^2_{X,Y}B-\nabla^2_{Y,X}B)
  (\omega^1,\ldots,\omega^k,V_1,\ldots,V_l)\notag\\
&\quad=\sum_{a=1}^k
B(\omega^1,\ldots,R(X,Y)^*\omega^a,\ldots,\omega^k,
  V_1,\ldots,V_l)\notag\\
&\qquad-\sum_{b=1}^l
B(\omega^1,\ldots,\omega^k,
  V_1,\ldots,R(X,Y)V_b,\ldots,V_l).
\tag{7.19}
\end{align}
In every local frame, these become
\begin{equation}
Z^i_{;pq}-Z^i_{;qp}=-R_{pqm}^iZ^m,
\tag{7.20}
\end{equation}
\begin{equation}
\beta_{j;pq}-\beta_{j;qp}=R_{pqj}^m\beta_m,
\tag{7.21}
\end{equation}
and
\begin{align}
B^{i_1\cdots i_k}_{j_1\cdots j_l;pq}
-B^{i_1\cdots i_k}_{j_1\cdots j_l;qp}
&=-R_{pqm}^{i_1}B^{m i_2\cdots i_k}_{j_1\cdots j_l}
-\cdots
-R_{pqm}^{i_k}B^{i_1\cdots i_{k-1}m}_{j_1\cdots j_l}
\notag\\
&\quad+R_{pqj_1}^mB^{i_1\cdots i_k}_{m j_2\cdots j_l}
+\cdots
+R_{pqj_l}^mB^{i_1\cdots i_k}_{j_1\cdots j_{l-1}m}.
\tag{7.22}
\end{align}
\end{theorem}

\begin{proof}
\sketch Symmetry of the connection and the formula for second total
covariant derivatives give, for every tensor field $B$,
\begin{align}
\nabla^2_{X,Y}B-\nabla^2_{Y,X}B
&=\nabla_X\nabla_YB-\nabla_Y\nabla_XB-\nabla_{[X,Y]}B.
\tag{7.23}
\end{align}
For $B=Z$, this is the definition of $R(X,Y)Z$, proving (7.17).

For a $1$-form $\beta$, repeated use of the covector derivative rule gives
\begin{align}
(\nabla_X\nabla_Y\beta)(Z)
&=XY(\beta(Z))-(\nabla_X\beta)(\nabla_YZ)
-\beta(\nabla_X\nabla_YZ)
-(\nabla_Y\beta)(\nabla_XZ),
\tag{7.24}\\
(\nabla_Y\nabla_X\beta)(Z)
&=YX(\beta(Z))-(\nabla_Y\beta)(\nabla_XZ)
-\beta(\nabla_Y\nabla_XZ)
-(\nabla_X\beta)(\nabla_YZ),
\tag{7.25}\\
(\nabla_{[X,Y]}\beta)(Z)
&=[X,Y](\beta(Z))-\beta(\nabla_{[X,Y]}Z).
\tag{7.26}
\end{align}
Substituting these into (7.23) leaves
$-\beta(R(X,Y)Z)$, proving (7.18).

The operator in (7.23) is a derivation for tensor products:
\begin{align*}
(\nabla^2_{X,Y}-\nabla^2_{Y,X})(F\otimes G)
&=(\nabla^2_{X,Y}F-\nabla^2_{Y,X}F)\otimes G\\
&\quad+F\otimes
(\nabla^2_{X,Y}G-\nabla^2_{Y,X}G).
\end{align*}
Applying this identity to decomposable tensors and using (7.17)--(7.18)
proves (7.19). For a local frame $(E_j)$ with dual coframe $(e^i)$, evaluate
at $(e^{i_1},\ldots,e^{i_k},E_{j_1},\ldots,E_{j_l})$ and use
\[
R(E_q,E_p)E_j=-R_{pqj}^mE_m,\qquad
R(E_q,E_p)^*e^i=-R_{pqm}^ie^m
\]
to obtain (7.22); (7.20) and (7.21) are its special cases.
\end{proof}

\section{Ricci and Scalar Curvatures}

On an $n$-dimensional Riemannian or pseudo-Riemannian manifold, the
\emph{Ricci tensor} is
\[
\mathrm{Rc}(X,Y)=\operatorname{tr}(Z\mapsto R(Z,X)Y).
\]
Its components and the \emph{scalar curvature} are
\[
R_{ij}=R_{kij}^{\,k}=g^{km}R_{kijm},
\qquad
S=\operatorname{tr}_g\mathrm{Rc}=R_i^i=g^{ij}R_{ij}.
\]

\begin{lemma}
\label{lem:ricci-symmetries}
\dcref{ch7:7.15}

The Ricci tensor is symmetric, and
\[
R_{ij}=R_{kij}^{\,k}=R^k{}_{ikj}
=-R^k{}_{kij}=-R_{ikj}^{\,k}.
\]
\end{lemma}
\begin{proof}
\sketch
Contract the first and fourth indices in the skew symmetries and pair
symmetry of the Riemann tensor. The first Bianchi identity then identifies
the two possible contractions and gives symmetry in the remaining indices.
\end{proof}

Define the \emph{traceless Ricci tensor} by
\[
\stackrel{\circ}{\mathrm{Rc}}
=\mathrm{Rc}-\frac1nSg.
\]

\begin{proposition}
\label{prop:ricci-tensor-decomposition}
\dcref{ch7:7.17}

Let $(M,g)$ be a Riemannian or pseudo-Riemannian $n$-manifold. Then
\[
\operatorname{tr}_g\stackrel{\circ}{\mathrm{Rc}}=0,
\]
and the orthogonal decomposition of the Ricci tensor is
\begin{equation}
\mathrm{Rc}=\stackrel{\circ}{\mathrm{Rc}}+\frac1nSg.
\tag{7.27}
\end{equation}
In the Riemannian case,
\begin{equation}
|\mathrm{Rc}|_g^2
=|\stackrel{\circ}{\mathrm{Rc}}|_g^2+\frac1nS^2.
\tag{7.28}
\end{equation}
\end{proposition}

\begin{remark}
The norm identity (7.28) uses positive definiteness. For a
pseudo-Riemannian metric, the corresponding identity is
\[
\langle\mathrm{Rc},\mathrm{Rc}\rangle_g
=\left\langle\stackrel{\circ}{\mathrm{Rc}},
\stackrel{\circ}{\mathrm{Rc}}\right\rangle_g+\frac1nS^2.
\]
\end{remark}

\begin{proof}
\sketch Since
$\operatorname{tr}_gg=g_{ij}g^{ji}=n$, the definition gives
$\operatorname{tr}_g\stackrel{\circ}{\mathrm{Rc}}=0$ and (7.27). For every
symmetric $2$-tensor $h$,
\[
\langle h,g\rangle_g
=g^{ik}g^{jl}h_{ij}g_{kl}
=g^{ij}h_{ij}
=\operatorname{tr}_gh.
\]
Thus $\stackrel{\circ}{\mathrm{Rc}}$ is orthogonal to $g$. In the Riemannian
case, take squared norms in (7.27) and use
$\langle g,g\rangle_g=n$ to obtain (7.28).
\end{proof}

For a smooth $2$-tensor $T$, define its \emph{exterior covariant derivative}
by
\[
(DT)(X,Y,Z)=-(\nabla T)(X,Y,Z)+(\nabla T)(X,Z,Y).
\qquad (7.29)
\]
In components,
\[
(DT)_{ijk}=-T_{ij;k}+T_{ik;j}.
\]

\begin{proposition}[Contracted Bianchi Identities]
\label{prop:contracted-bianchi-identities}
\uses{lem:ricci-symmetries, prop:connection-on-tensor-bundles, prop:musical-isomorphisms-commute, prop:exterior-derivative-covariant-formula, prob:exterior-derivative-covariant-formula}
\dcref{ch7:7.18}


On a Riemannian or pseudo-Riemannian manifold,
\[
\operatorname{tr}_g(\nabla Rm)=-D(\mathrm{Rc}),
\qquad (7.30)
\]
and
\[
\operatorname{tr}_g(\nabla\mathrm{Rc})=\tfrac12dS,
\qquad (7.31)
\]
where each trace contracts the first and last indices. In components,
\[
R_{ijkl;i}=R_{jk;l}-R_{jl;k},
\qquad (7.32)
\]
and
\[
R_{il;i}=\tfrac12S_{;l}.
\qquad (7.33)
\]
\end{proposition}

\begin{proof}
\sketch Raise the index $m$ in (7.15) and contract $i,m$. Covariant
differentiation commutes with raising indices and contraction, and the
curvature symmetries reduce the result to (7.32). Contract (7.32) in $j,k$
and simplify to obtain (7.33). Expanding the contractions and $D$ gives
(7.30)--(7.31).
\end{proof}

For the curvature convention used here, Ricci curvature is the trace of the
curvature endomorphism in its first and last indices. Under the opposite
curvature sign convention, it must instead be defined by tracing $Rm$ in its
first and third, equivalently second and fourth, indices. Traces in the first
two or last two indices vanish by antisymmetry.

A Riemannian or pseudo-Riemannian metric is \emph{Einstein} if
\begin{equation}
\mathrm{Rc}=\lambda g
\qquad\text{for some constant }\lambda.
\tag{7.34}
\end{equation}

\begin{proposition}[Schur's Lemma]
\label{prop:schur-lemma}
\uses{prop:contracted-bianchi-identities}
\dcref{ch7:7.19}


Let $(M,g)$ be a connected Riemannian or pseudo-Riemannian manifold of
dimension $n\geq3$. If $\mathrm{Rc}=fg$ for a smooth function
$f:M\to\mathbb R$, then $f$ is constant and $g$ is Einstein.
\end{proposition}

\begin{proof}
\sketch Taking the trace of $\mathrm{Rc}=fg$ gives $f=S/n$, hence
$\stackrel{\circ}{\mathrm{Rc}}=0$ and
$\nabla\stackrel{\circ}{\mathrm{Rc}}=0$. In coordinates,
\[
0=R_{ij;k}-\frac1nS_{;k}g_{ij}.
\]
Contracting $i,k$ and applying (7.33) gives
\[
0=\left(\frac12-\frac1n\right)S_{;j}.
\]
Since $n\geq3$, $dS=0$. Connectedness makes $S$, and therefore $f$, constant.
\end{proof}

\begin{corollary}
\label{cor:einstein-zero-ricci}
\uses{prop:schur-lemma}
\dcref{ch7:7.20}


If $(M,g)$ is a connected Riemannian or pseudo-Riemannian manifold of
dimension $n\geq3$, then $g$ is Einstein if and only if
\[
\stackrel{\circ}{\mathrm{Rc}}=0.
\]
\end{corollary}

\begin{proof}
\sketch If $\mathrm{Rc}=\lambda g$, then tracing gives $\lambda=S/n$, so
$\stackrel{\circ}{\mathrm{Rc}}=0$. Conversely, vanishing traceless Ricci gives
$\mathrm{Rc}=(S/n)g$, and Schur's lemma makes $S/n$ constant.
\end{proof}

On a Lorentzian $4$-manifold, the Einstein field equation has the form
\begin{equation}
\mathrm{Rc}-\tfrac12Sg=T,
\tag{7.35}
\end{equation}
where $T$ is a symmetric $2$-tensor. In vacuum, $T=0$; tracing (7.35) in
dimension $4$ gives $S=2S$, hence $S=0$ and $\mathrm{Rc}=0$. With a
cosmological term $\Lambda g$ on the left, the vacuum equation instead has
the form $\mathrm{Rc}=\lambda g$ for a constant $\lambda$.
\section{The Weyl Tensor}

Let $V$ be an $n$-dimensional real vector space.  An \emph{algebraic
curvature tensor} is a covariant $4$-tensor $T$ satisfying
\begin{enumerate}
\item[(a)] $T(w,x,y,z)=-T(x,w,y,z)$;
\item[(b)] $T(w,x,y,z)=-T(w,x,z,y)$;
\item[(c)] $T(w,x,y,z)=T(y,z,w,x)$;
\item[(d)] $T(w,x,y,z)+T(x,y,w,z)+T(y,w,x,z)=0$.
\end{enumerate}
Write $\mathscr R(V^*)\subseteq T^4(V^*)$ for their vector space.  Although
(c) follows from (a), (b), and (d), all four conditions are included in this
notation.

\begin{proposition}[Dimension of the Algebraic Curvature Tensors]
\label{prop:dim-algebraic-curvature-tensors}
\uses{prop:curvature-tensor-symmetries}
\dcref{ch7:7.21}

If $\dim V=n$, then
\begin{equation}
\dim \mathscr{R}(V^*) = \frac{n^2(n^2-1)}{12}. \tag{7.36}
\end{equation}
\end{proposition}

\begin{proof}
\sketch
Let $\mathscr B(V^*)$ be the tensors satisfying (a)--(c).  The map
\[
\Phi:\Sigma^2((\Lambda^2V)^*)\longrightarrow\mathscr B(V^*),
\qquad
\Phi(B)(w,x,y,z)=B(w\wedge x,y\wedge z),
\]
is an isomorphism.  Hence
\[
\dim\mathscr B(V^*)
=\frac{\binom n2(\binom n2+1)}2
=\frac{n(n-1)(n^2-n+2)}8.
\]
For $T\in\mathscr B(V^*)$, set
\[
\pi(T)(w,x,y,z)=\frac13\bigl(T(w,x,y,z)+T(x,y,w,z)+T(y,w,x,z)\bigr).
\]
The curvature tensors form $\ker\pi$.  On $\mathscr B(V^*)$, the map $\pi$
is alternation and maps onto $\Lambda^4(V^*)$, because every $4$-form lies
in $\mathscr B(V^*)$ and is fixed by alternation.  Rank--nullity gives
\[
\dim\mathscr R(V^*)
=\frac{n(n-1)(n^2-n+2)}8-\binom n4,
\]
which simplifies to (7.36).
\end{proof}

Now give $V$ a possibly indefinite scalar product $g$.  For symmetric
$2$-tensors $h,k$, their \emph{Kulkarni--Nomizu product} is
\begin{equation}
h \owedge k(w,x,y,z) = h(w,z)k(x,y) + h(x,y)k(w,z)
{}- h(w,y)k(x,z) - h(x,z)k(w,y). \tag{7.37}
\end{equation}
Its components are
\[
(h\owedge k)_{ijlm}
=h_{im}k_{jl}+h_{jl}k_{im}-h_{il}k_{jm}-h_{jm}k_{il}.
\]
This sign agrees with the curvature convention used here.

\begin{lemma}[Properties of the Kulkarni--Nomizu Product]
\label{lem:kulkarni-nomizu-properties}
\dcref{ch7:7.22}
Let $(V,g)$ be an $n$-dimensional scalar-product space, let $h,k$ be
symmetric $2$-tensors, let $T$ be an algebraic curvature tensor, and trace on
the first and last indices.
\begin{enumerate}
\item[(a)] $h\owedge k$ is an algebraic curvature tensor.
\item[(b)] $h\owedge k=k\owedge h$.
\item[(c)] $\operatorname{tr}_g(h\owedge g)
=(n-2)h+(\operatorname{tr}_gh)g$.
\item[(d)] $\operatorname{tr}_g(g\owedge g)=2(n-1)g$.
\item[(e)] $\langle T,h\owedge g\rangle_g
=4\langle\operatorname{tr}_gT,h\rangle_g$.
\item[(f)] If $g$ is positive definite, then
$|g\owedge h|_g^2=4(n-2)|h|_g^2+4(\operatorname{tr}_gh)^2$.
\end{enumerate}
\end{lemma}

\begin{proof}
\sketch
Formula (7.37) immediately gives antisymmetry in each pair, pair symmetry,
and commutativity in $h,k$.  The cyclic sum in the first three arguments
vanishes by symmetry of $h,k$, proving the algebraic Bianchi identity.  In a
basis,
\begin{align*}
(\operatorname{tr}_g(h\owedge g))_{jl}
&=g^{im}(h_{im}g_{jl}+h_{jl}g_{im}-h_{il}g_{jm}-h_{jm}g_{il})\\
&=(\operatorname{tr}_gh)g_{jl}+(n-2)h_{jl},
\end{align*}
which proves (c), and (d) follows by taking $h=g$.  The remaining identities
follow by the corresponding contractions of (7.37).
\end{proof}

\begin{proposition}[A Right Inverse to Curvature Trace]
\label{prop:right-inverse-trace}
\uses{lem:kulkarni-nomizu-properties}
\dcref{ch7:7.23}

Let $(V,g)$ be an $n$-dimensional scalar-product space with $n\geq3$, and
define
\begin{equation}
G(h) = \frac{1}{n-2}\left(h - \frac{\operatorname{tr}_g h}{2(n-1)}g\right)
\owedge g. \tag{7.38}
\end{equation}
Then $G:\Sigma^2(V^*)\to\mathscr R(V^*)$ is a right inverse of
$\operatorname{tr}_g$, and
$\operatorname{Im}G=\ker(\operatorname{tr}_g)^\perp$.
\end{proposition}

\begin{proof}
\sketch
Parts (c) and (d) of the preceding lemma give
$\operatorname{tr}_gG(h)=h$.  Thus $G$ is injective and the trace is
surjective, so $\operatorname{Im}G$ and
$\ker(\operatorname{tr}_g)^\perp$ have the same dimension.  If
$T\in\ker(\operatorname{tr}_g)$, part (e) gives
$\langle T,G(h)\rangle_g=0$ for every $h$, proving equality.
\end{proof}

For a Riemannian or pseudo-Riemannian metric of dimension $n\geq3$, define
the \emph{Schouten tensor} and \emph{Weyl tensor} by
\[
P=\frac1{n-2}\left(\mathrm{Rc}-\frac{S}{2(n-1)}g\right)
\]
and
\[
W=\mathrm{Rm}-P\owedge g
=\mathrm{Rm}-\frac1{n-2}\mathrm{Rc}\owedge g
+\frac{S}{2(n-1)(n-2)}g\owedge g.
\]

\begin{proposition}[Weyl Decomposition]
\label{prop:weyl-decomposition}
\uses{prop:right-inverse-trace}
\dcref{ch7:7.24}

For every Riemannian or pseudo-Riemannian $n$-manifold with $n\geq3$,
$\operatorname{tr}_gW=0$, and
\[
\mathrm{Rm}=W+P\owedge g
\]
is the orthogonal decomposition associated with
$\mathscr R(V^*)=\ker(\operatorname{tr}_g)
\oplus\ker(\operatorname{tr}_g)^\perp$.
\end{proposition}

\begin{proof}
\sketch
Formula (7.38) gives
$P\owedge g=G(\mathrm{Rc})=G(\operatorname{tr}_g\mathrm{Rm})$; the result is
the trace decomposition from the preceding proposition.
\end{proof}

\begin{corollary}[Low-Dimensional Curvature Spaces]
\label{cor:curvature-space-dimensions}
\uses{prop:dim-algebraic-curvature-tensors, lem:kulkarni-nomizu-properties, prop:right-inverse-trace}
\dcref{ch7:7.25}

Let $(V,g)$ be an $n$-dimensional real scalar-product space.
\begin{enumerate}
\item[(a)] If $n=0$ or $n=1$, then $\mathscr R(V^*)=\{0\}$.
\item[(b)] If $n=2$, then $\mathscr R(V^*)$ is one-dimensional and spanned by
$g\owedge g$.
\item[(c)] If $n=3$, then $\mathscr R(V^*)$ is six-dimensional and
$G:\Sigma^2(V^*)\to\mathscr R(V^*)$ is an isomorphism.
\end{enumerate}
\end{corollary}

\begin{proof}
\sketch
Formula (7.36) gives the dimensions.  For $n=2$,
$\operatorname{tr}_g(g\owedge g)=2g\ne0$, so $g\owedge g$ spans.  For $n=3$,
$G$ is injective and both its domain and codomain have dimension six.
\end{proof}

\begin{corollary}[The Curvature Tensor in Dimension 3]
\label{cor:ricci-determines-curvature-dim3}
\uses{cor:curvature-space-dimensions, prop:weyl-decomposition}
\dcref{ch7:7.26}

On every Riemannian or pseudo-Riemannian $3$-manifold, $W=0$ and
\begin{equation}
\mathrm{Rm}=P\owedge g
=\mathrm{Rc}\owedge g-\frac14S\,g\owedge g. \tag{7.39}
\end{equation}
Thus the Ricci tensor determines the full curvature tensor.
\end{corollary}

\begin{proof}
\sketch
In dimension three, $G$ and $\operatorname{tr}_g$ are inverse isomorphisms.
Since $\operatorname{tr}_gW=0$, one has $W=0$; the displayed formula follows
from the definition of $P$.
\end{proof}

\begin{corollary}[The Curvature Tensor in Dimension 2]
\label{cor:curvature-tensor-dim-2}
\uses{cor:curvature-space-dimensions, lem:kulkarni-nomizu-properties}
\dcref{ch7:7.27}

On every Riemannian or pseudo-Riemannian $2$-manifold,
\[
\mathrm{Rm}=\frac14Sg\owedge g,
\qquad
\mathrm{Rc}=\frac12Sg.
\]
\end{corollary}

\begin{proof}
\sketch
The one-dimensional result gives $\mathrm{Rm}=f,g\owedge g$.  Tracing twice
yields $\mathrm{Rc}=2fg$ and $S=4f$, proving both formulas.
\end{proof}

On a $2$-manifold, the traceless Ricci tensor always vanishes, but $S$ need
not be constant; the metric is Einstein exactly when $S$ is constant.

\begin{proposition}[The Ricci Decomposition of the Curvature Tensor]
\label{prop:ricci-decomposition}
\uses{cor:curvature-tensor-dim-2, lem:kulkarni-nomizu-properties, prop:ricci-tensor-decomposition}
\dcref{ch7:7.28}

Let $(M,g)$ be a Riemannian or pseudo-Riemannian $n$-manifold with $n\geq3$.
Then
\begin{equation}
\mathrm{Rm}=W+\frac1{n-2}\stackrel{\circ}{\mathrm{Rc}}\owedge g
+\frac1{2n(n-1)}Sg\owedge g. \tag{7.40}
\end{equation}
This decomposition is orthogonal.  In the Riemannian case,
\[
|\mathrm{Rm}|_g^2=|W|_g^2
+\frac1{(n-2)^2}|\stackrel{\circ}{\mathrm{Rc}}\owedge g|_g^2
+\frac1{4n^2(n-1)^2}|Sg\owedge g|_g^2
\]
\begin{equation}
=|W|_g^2+\frac4{n-2}|\stackrel{\circ}{\mathrm{Rc}}|_g^2
+\frac2{n(n-1)}S^2. \tag{7.41}
\end{equation}
\end{proposition}

\begin{proof}
\sketch
Substitute the decomposition of $\mathrm{Rc}$ into the Schouten tensor.
Part (e) of the Kulkarni--Nomizu lemma proves orthogonality, and part (f)
gives the coefficients in (7.41).
\end{proof}

\section{Curvatures of Conformally Related Metrics}

Metrics $g$ and $\widetilde g$ are \emph{conformal} if
$\widetilde g=e^{2f}g$ for some smooth real-valued $f$.

\begin{proposition}[Conformal Transformation of the Levi-Civita Connection]
\label{prop:conformal-transformation-levi-civita}
\uses{thm:fundamental-theorem-riemannian-geometry}
\dcref{ch7:7.29}

Let $(M,g)$ be a Riemannian or pseudo-Riemannian $n$-manifold, with or without
boundary, and let $\widetilde g=e^{2f}g$.  Their Levi-Civita connections satisfy
\[
\widetilde\nabla_XY=\nabla_XY+(Xf)Y+(Yf)X
-\langle X,Y\rangle_g\operatorname{grad}f, \quad (7.42)
\]
and, in local coordinates,
\[
\widetilde\Gamma^k_{ij}=\Gamma^k_{ij}+f_i\delta^k_j+f_j\delta^k_i
-g^{kl}f_lg_{ij}, \quad (7.43)
\]
where $f_i=\partial_if$ and the gradient is computed using $g$.
\end{proposition}

\begin{proof}
\sketch
The connection defined by the right side of (7.42) is symmetric.  Direct
differentiation of $\widetilde g=e^{2f}g$ shows that it is compatible with
$\widetilde g$, so uniqueness of the Levi-Civita connection proves (7.42).
Expansion in a coordinate frame gives (7.43).
\end{proof}

\begin{theorem}[Conformal Transformation of the Curvature]
\label{thm:conformal-transformation-curvature}
\uses{prop:conformal-transformation-levi-civita}
\dcref{ch7:7.30}

Let $g$ be a Riemannian or pseudo-Riemannian metric on an $n$-manifold with or
without boundary, let $f\in C^\infty(M)$, and set $\widetilde g=e^{2f}g$.
In the Riemannian case,
\[
\widetilde{\mathrm{Rm}}=e^{2f}\left(\mathrm{Rm}-(\nabla^2f)\owedge g
+(df\otimes df)\owedge g-\tfrac12|df|_g^2(g\owedge g)\right), \quad (7.44)
\]
\[
\widetilde{\mathrm{Rc}}=\mathrm{Rc}-(n-2)\nabla^2f
+(n-2)(df\otimes df)-(\Delta f+(n-2)|df|_g^2)g, \quad (7.45)
\]
\[
\widetilde S=e^{-2f}\left(S-2(n-1)\Delta f
-(n-1)(n-2)|df|_g^2\right), \quad (7.46)
\]
where all quantities on the right are computed using $g$ and
$\Delta f=\operatorname{div}(\operatorname{grad}f)$.  If $n\geq3$, then
\[
\widetilde P=P-\nabla^2f+(df\otimes df)-\tfrac12|df|_g^2g, \quad (7.47)
\]
\[
\widetilde W=e^{2f}W. \quad (7.48)
\]
For pseudo-Riemannian $g$, the same formulas hold with every $|df|_g^2$
replaced by $\langle df,df\rangle_g$.
\end{theorem}

\begin{proof}
\sketch
At an arbitrary point choose $g$-normal coordinates and substitute (7.43)
into the coordinate curvature formula.  Since
$g_{ij}=\delta_{ij}$, $\partial_kg_{ij}=0$, and $\Gamma^k_{ij}=0$ at that
point, simplification gives
\begin{align*}
\widetilde R_{ijkl}=e^{2f}\bigl(&R_{ijkl}
-(f_{;il}g_{jk}+f_{;jk}g_{il}-f_{;ik}g_{jl}-f_{;jl}g_{ik})\\
&+(f_{;i}f_{;l}g_{jk}+f_{;j}f_{;k}g_{il}
-f_{;i}f_{;k}g_{jl}-f_{;j}f_{;l}g_{ik})\\
&-g^{mp}f_{;m}f_{;p}(g_{il}g_{jk}-g_{ik}g_{jl})\bigr),
\end{align*}
which is (7.44).  Trace with
$\widetilde g^{ij}=e^{-2f}g^{ij}$ and use the Kulkarni--Nomizu trace
identities to obtain (7.45) and (7.46).  Substitution into the definition of
$P$ gives (7.47), and subtracting
$\widetilde P\owedge\widetilde g$ from $\widetilde{\mathrm{Rm}}$ gives (7.48).
\end{proof}

A metric is \emph{locally conformally flat} if every point has a neighborhood
conformally equivalent to an open subset of Euclidean space, or of the
corresponding pseudo-Euclidean space in the pseudo-Riemannian case.

\begin{corollary}[Conformal Flatness Forces Vanishing Weyl Tensor]
\label{cor:weyl-zero-locally-conformally-flat}
\uses{thm:conformal-transformation-curvature}
\dcref{ch7:7.31}

If a Riemannian or pseudo-Riemannian manifold of dimension $n\geq3$ is locally
conformally flat, then its Weyl tensor vanishes identically.
\end{corollary}

\begin{proof}
\sketch
Locally, a conformal multiple $\widetilde g=e^{2f}g$ is flat and therefore has
$\widetilde W=0$.  Formula (7.48) implies $W=0$.
\end{proof}

For $n\geq3$, define the \emph{Cotton tensor} by $C=-DP$; in a local frame,
\begin{equation}
C_{ijk}=P_{ij;k}-P_{ik;j}. \tag{7.49}
\end{equation}

\begin{proposition}[Divergence of the Weyl Tensor]
\label{prop:trace-covariant-derivative-weyl}
\uses{prop:contracted-bianchi-identities}
\dcref{ch7:7.32}

On every Riemannian or pseudo-Riemannian $n$-manifold with $n\geq3$,
\begin{equation}
\operatorname{tr}_g(\nabla W)=(n-3)C, \tag{7.50}
\end{equation}
where the trace is on the first and last indices of $\nabla W$.
\end{proposition}

\begin{proof}
\sketch
Differentiate $W=\mathrm{Rm}-P\owedge g$ and use the first contracted Bianchi
identity to obtain
\begin{equation}
W_{ijkl;i}=R_{jk;l}-R_{jl;k}-P_{il;i}g_{jk}-P_{jk;i}g_{il}
+P_{ik;i}g_{jl}+P_{jl;i}g_{ik}. \tag{7.51}
\end{equation}
The second contracted Bianchi identity and the definition of $P$ give
\[
P_{il;i}=\frac1{2(n-1)}S_{;l},
\qquad
P_{ik;i}=\frac1{2(n-1)}S_{;k}.
\]
Substitution into (7.51) yields
\[
W_{ijkl;i}=(n-3)(P_{jk;l}-P_{jl;k})=(n-3)C_{jkl}.
\]
\end{proof}

\begin{corollary}[Vanishing Weyl Tensor Forces Vanishing Cotton Tensor]
\label{cor:weyl-cotton-vanish-dimension-4}
\uses{prop:trace-covariant-derivative-weyl}
\dcref{ch7:7.33}

If $\dim M\geq4$ and the Weyl tensor of a Riemannian or pseudo-Riemannian
manifold vanishes identically, then its Cotton tensor also vanishes.
\end{corollary}

\begin{proposition}[Conformal Invariance of the Cotton Tensor in Dimension 3]
\label{prop:conformal-invariance-cotton-dim3}
\dcref{ch7:7.34}
If $(M,g)$ is a Riemannian or pseudo-Riemannian $3$-manifold and
$\widetilde g=e^{2f}g$, then $\widetilde C=C$.
\end{proposition}

\begin{corollary}[Conformal Flatness Forces Vanishing Cotton Tensor]
\label{cor:cotton-tensor-vanishes-conformally-flat}
\uses{prop:conformal-invariance-cotton-dim3}
\dcref{ch7:7.35}

The Cotton tensor of every locally conformally flat $3$-manifold vanishes
identically.
\end{corollary}
\begin{proof}
\sketch
In conformally flat coordinates the metric is conformal to a flat metric,
whose Cotton tensor vanishes. Conformal invariance in dimension three then
gives the result.
\end{proof}

\begin{theorem}[Weyl--Schouten]
\label{thm:weyl-schouten}
\uses{cor:weyl-zero-locally-conformally-flat, cor:cotton-tensor-vanishes-conformally-flat, cor:ricci-determines-curvature-dim3, cor:weyl-cotton-vanish-dimension-4, lem:overdetermined-covector-system}
\dcref{ch7:7.37}

Let $(M,g)$ be a Riemannian or pseudo-Riemannian $n$-manifold with $n\geq3$.
\begin{enumerate}
\item[(a)] If $n\geq4$, then $(M,g)$ is locally conformally flat if and only
if $W=0$.
\item[(b)] If $n=3$, then $(M,g)$ is locally conformally flat if and only if
$C=0$.
\end{enumerate}
\end{theorem}

\begin{proof}
\sketch
Necessity follows from the two preceding conformal-flatness corollaries.  For
sufficiency, the relevant hypothesis and the dimension-three or
dimension-at-least-four consequences imply $W=C=0$.  Since
$\widetilde W=e^{2f}W$, every $\widetilde g=e^{2f}g$ also has zero Weyl tensor,
and it is flat once $\widetilde P=0$.  Formula (7.47) shows that this condition
is
\begin{equation}
P-\nabla^2f+(df\otimes df)
-\tfrac12\langle df,df\rangle_gg=0. \tag{7.52}
\end{equation}

For a covector $\xi$, define
\[
A(\xi)=(\xi\otimes\xi)-\tfrac12\langle\xi,\xi\rangle_gg+P,
\]
or
\begin{equation}
A(\xi)_{ij}=\xi_i\xi_j-\tfrac12\xi_m\xi^mg_{ij}+P_{ij}. \tag{7.53}
\end{equation}
Equation (7.52) becomes $\nabla(df)=A(df)$.  The Ricci identity gives the
necessary compatibility relation
\begin{equation}
A(\xi)_{ij;k}-A(\xi)_{ik;j}=R_{jki}^{l}\xi_l. \tag{7.54}
\end{equation}
After differentiating (7.53), replacing every $\nabla\xi$ by $A(\xi)$, and
using $\mathrm{Rm}=W+P\owedge g$, the obstruction is
\[
A(\xi)_{ij;k}-A(\xi)_{ik;j}-R_{jki}^{l}\xi_l
=-W_{jki}^{l}\xi_l+C_{ijk}. \quad (7.55)
\]
It vanishes because $W=C=0$.  The following lemma therefore produces a local
solution of $\nabla\xi=A(\xi)$.  Since $A(\xi)$ is symmetric, $\xi$ is closed;
the Poincar\'e lemma gives $\xi=df$ locally, and (7.52) makes the conformal
metric flat.
\end{proof}

\begin{lemma}[Local Solvability of the Covector System]
\label{lem:overdetermined-covector-system}
\dcref{ch7:7.38}
Let $(M,g)$ be Riemannian or pseudo-Riemannian, and consider
\[
\nabla\xi=A(\xi), \quad (7.56)
\]
for a smooth map $A:T^*M\to T^2T^*M$.  Assume that, for every smooth covector
field $\xi$, replacing each occurrence of $\nabla\xi$ in
$\nabla(A(\xi))$ by $A(\xi)$ yields
\[
A(\xi)_{ij;k}-A(\xi)_{ik;j}=R_{jki}^{l}\xi_l. \quad (7.57)
\]
Then for every $p\in M$ and $\eta_0\in T_p^*M$, a smooth local solution exists
with $\xi_p=\eta_0$.
\end{lemma}

\begin{proof}
\sketch
In local coordinates, the system is
\[
\frac{\partial\xi_i}{\partial x^j}=a_{ij}(x,\xi),
\qquad
a_{ij}(x,\xi)=\Gamma^k_{ij}(x)\xi_k+A(\xi)_{ij}.
\]
The Frobenius theorem [LeeSM, Prop.~19.29] gives a local solution with
arbitrary initial value if
\begin{equation}
\frac{\partial a_{ij}}{\partial x^k}
+a_{kl}\frac{\partial a_{ij}}{\partial\xi_l}
=\frac{\partial a_{ik}}{\partial x^j}
+a_{jl}\frac{\partial a_{ik}}{\partial\xi_l}. \tag{7.58}
\end{equation}
By the chain rule, this is symmetry in $j,k$ after replacing
$\partial_k\xi_l$ by $a_{lk}$.  Since
$\xi_{l;k}=\partial_k\xi_l-\Gamma^m_{lk}\xi_m$, that replacement is exactly
$\nabla\xi=A(\xi)$, and expansion of the assumed compatibility identity is
(7.58).
\end{proof}

Every Riemannian $1$-manifold is flat.  Every Riemannian or
pseudo-Riemannian $2$-manifold is locally conformally flat; the corresponding
local conformal charts are called \emph{isothermal coordinates}.

\section{Further Curvature Results}

\begin{proposition}[Locally symmetric spaces have parallel curvature]
\label{prop:symmetric-space-parallel-curvature}
\dcref{ch7:7-3}
Every Riemannian locally symmetric space has parallel curvature tensor:
\[
\nabla\operatorname{Rm}=0.
\]
\end{proposition}

\begin{exercise}
\label{prob:symmetric-space-parallel-curvature}
\dcref{ch7:7-3}
Show that the curvature tensor of a Riemannian locally symmetric space is parallel: $\nabla \text{Rm} = 0$. (Used on pp. 297, 351.)
\end{exercise}

\begin{proof}
\sketch
At each point $p$, the local geodesic symmetry is an isometry fixing $p$ with
differential $-\operatorname{Id}$. It preserves $\nabla\operatorname{Rm}$ by
local-isometry invariance, while its differential acts by $-1$ on this
covariant $5$-tensor. Hence $(\nabla\operatorname{Rm})_p=0$.
\end{proof}

\begin{proposition}[Curvature of a product metric]
\label{prop:product-metric-curvature}
\dcref{ch7:7-6}
\uses{lem:pullback-immersion-metric}
Let $(M_1,g_1)$ and $(M_2,g_2)$ be Riemannian manifolds, give
$M_1\times M_2$ the product metric $g=g_1\oplus g_2$, and write
$\pi_i\colon M_1\times M_2\to M_i$ for the projections. Then
\begin{align*}
\operatorname{Rm}&=\pi_1^*\operatorname{Rm}_1+\pi_2^*\operatorname{Rm}_2,\\
\operatorname{Rc}&=\pi_1^*\operatorname{Rc}_1+\pi_2^*\operatorname{Rc}_2,\\
S&=\pi_1^*S_1+\pi_2^*S_2.
\end{align*}
\end{proposition}

\begin{exercise}
\label{prob:product-metric-curvature}
\dcref{ch7:7-6}
\uses{exer:prove-lemma-2-11}
Suppose $(M_1, g_1)$ and $(M_2, g_2)$ are Riemannian manifolds, and $M_1 \times M_2$ is endowed with the product metric $g = g_1 \oplus g_2$ as in (2.12). Show that the Riemann curvature, Ricci curvature, and scalar curvature of $g$ are given by the following formulas:
\begin{align*}
Rm &= \pi_1^* Rm_1 + \pi_2^* Rm_2, \\
Rc &= \pi_1^* Rc_1 + \pi_2^* Rc_2, \\
S &= \pi_1^* S_1 + \pi_2^* S_2,
\end{align*}
where $Rm_i$, $Rc_i$, and $S_i$ are the Riemann, Ricci, and scalar curvatures of $(M_i, g_i)$, and $\pi_i : M_1 \times M_2 \to M_i$ is the projection. (Used on pp. 257, 261.)
\end{exercise}

\begin{proof}
\sketch
The Levi--Civita connection of the product metric splits into the pullbacks of
the two factor connections. Its curvature therefore has no mixed terms and
is the sum of the pulled-back factor curvatures. Taking the metric traces gives
the Ricci and scalar-curvature identities.
\end{proof}

\begin{proposition}[O'Neill's curvature formula]
\label{prop:oneill-formula}
\dcref{ch7:7-14}
\uses{def:riemannian-submersion,prop:horizontal-lift-of-vector-field}
Let $\pi:(\widetilde M,\widetilde g)\to(M,g)$ be a Riemannian submersion,
let $W,X,Y,Z$ be vector fields on $M$, and let
$\widetilde W,\widetilde X,\widetilde Y,\widetilde Z$ be their horizontal
lifts. Set
$A_{UV}=[\widetilde U,\widetilde V]^V$ and
$\widetilde R=\widetilde{\operatorname{Rm}}(
\widetilde W,\widetilde X,\widetilde Y,\widetilde Z)$. Then
\[
\begin{gathered}
\operatorname{Rm}(W,X,Y,Z)\circ\pi\\
=\widetilde R-\frac12\langle A_{WX},A_{YZ}\rangle\\
\quad-\frac14\langle A_{WY},A_{XZ}\rangle\\
\quad+\frac14\langle A_{WZ},A_{XY}\rangle.
\end{gathered}
\]
\end{proposition}

\begin{exercise}[O'Neill's formula]
\label{prob:oneill-formula}
\dcref{ch7:7-14}
\uses{prob:riemannian-submersion-connection}
Suppose $\pi : (\tilde{M}, \tilde{g}) \to (M, g)$ is a Riemannian submersion. Using the notation and results of Problem 5-6, prove O'Neill's formula:
\begin{align*}
\mathrm{Rm}(W, X, Y, Z) \circ \pi &= \widetilde{\mathrm{Rm}}(\tilde{W}, \tilde{X}, \tilde{Y}, \tilde{Z}) - \frac{1}{2}\langle[\tilde{W}, \tilde{X}]^V, [\tilde{Y}, \tilde{Z}]^V\rangle \\
&\quad - \frac{1}{4}\langle[\tilde{W}, \tilde{Y}]^V, [\tilde{X}, \tilde{Z}]^V\rangle + \frac{1}{4}\langle[\tilde{W}, \tilde{Z}]^V, [\tilde{X}, \tilde{Y}]^V\rangle.
\end{align*}
(Used on p. 258.)
\end{exercise}

\begin{proof}
\sketch Decompose the Levi--Civita derivatives of horizontal lifts into their
horizontal and vertical parts, substitute those decompositions into the
curvature definition, and use metric compatibility. The vertical bracket
terms combine with coefficients $-\frac12$, $-\frac14$, and $\frac14$.
\end{proof}

\begin{exercise}
\label{exer:constructed-vector-field-smooth}
\dcref{ch7:7.9}
Prove that the vector field $V$ constructed in the preceding proof is smooth.
\end{exercise}

\begin{exercise}
\label{exer:prove-cotton-tensor-vanishes}
\dcref{ch7:7.36}
\uses{cor:cotton-tensor-vanishes-conformally-flat}
Prove this corollary.
\end{exercise}

\begin{exercise}
\label{exer:prove-curvature-isometry}
\dcref{ch7:7.7}
\uses{prop:curvature-isometry-invariant}
Prove Proposition 7.6.
\end{exercise}

\begin{exercise}
\label{exer:ricci-symmetries-proof}
\dcref{ch7:7.16}
\uses{lem:ricci-symmetries}
Prove Lemma 7.15, using the symmetries of the curvature tensor.
\end{exercise}

\begin{exercise}[Bochner's formula]
\label{prob:bochner-formula}
\dcref{ch7:7-7}
\uses{exer:hilbert-axioms-euclidean-plane,thm:ricci-identities}
Suppose $(M, g)$ is a Riemannian manifold and $u \in C^\infty(M)$. Use (5.29) and (7.21) to prove \textbf{Bochner's formula}:
\[
\Delta\left(\tfrac{1}{2} |\operatorname{grad} u|^2\right) = |\nabla^2 u|^2 + \langle \operatorname{grad}(\Delta u), \operatorname{grad} u \rangle + Rc(\operatorname{grad} u, \operatorname{grad} u).
\]
\end{exercise}

\begin{exercise}
\label{prob:cartan-second-structure-equation}
\dcref{ch7:7-5}
\uses{prob:connection-1forms-cartan}
Let $\nabla$ be the Levi-Civita connection on a Riemannian or pseudo-Riemannian manifold $(M,g)$, and let $\omega_i^j$ be its connection 1-forms with respect to a local frame $(E_i)$ (Problem 4-14). Define a matrix of 2-forms $\Omega_i^j$, called the \textit{curvature 2-forms}, by
\[
\Omega_i^j = \frac{1}{2}R_{klij} e^k \wedge e^l,
\]
where $(e^i)$ is the coframe dual to $(E_i)$. Show that the curvature 2-forms satisfy \textit{Cartan's second structure equation}:
\[
\Omega_i^j = d\omega_i^j - \omega_k^j \wedge \omega_i^k.
\]
[Hint: Expand $R(E_k, E_l)E_i$ in terms of $\nabla$ and $\omega_i^j$.]
\end{exercise}

\begin{exercise}
\label{prob:conformal-laplacian-yamabe}
\dcref{ch7:7-11}
Let $(M, g)$ be a Riemannian manifold of dimension $n > 2$. Define the \textbf{conformal Laplacian} $L: C^\infty(M) \to C^\infty(M)$ by the formula
\[
Lu = -\frac{4(n-1)}{n-2} \Delta u + Su,
\]
where $\Delta$ is the Laplace--Beltrami operator of $g$ and $S$ is its scalar curvature. Prove that if $\tilde{g} = e^{2f} g$ for some $f \in C^\infty(M)$, and $\tilde{L}$ denotes the conformal Laplacian with respect to $\tilde{g}$, then for every $u \in C^\infty(M)$,
\[
e^{\frac{n+2}{2}f} \tilde{L}u = L\left(e^{\frac{n-2}{2}f} u\right).
\]
Conclude that a metric $\tilde{g}$ conformal to $g$ has constant scalar curvature $\lambda$ if and only if it can be expressed in the form $\tilde{g} = \varphi^{4/(n-2)} g$, where $\varphi$ is a smooth positive solution to the \textbf{Yamabe equation}:
\begin{equation}
L\varphi = \lambda \varphi^{\frac{n+2}{n-2}}.
\tag{7.59}
\end{equation}
\end{exercise}

\begin{exercise}
\label{prob:cotton-tensor-conformal}
\dcref{ch7:7-10}
\uses{prop:conformal-invariance-cotton-dim3}
Prove Proposition 7.34 (conformal invariance of the Cotton tensor in dimension 3).
\end{exercise}

\begin{exercise}
\label{prob:curvature-coordinate-formula}
\dcref{ch7:7-2}
\uses{prop:curvature-tensor-components}
Prove Proposition 7.4 (the formula for the curvature tensor in coordinates).
\end{exercise}

\begin{exercise}
\label{prob:curvature-multilinearity-in-functions}
\dcref{ch7:7-1}
\uses{prop:curvature-multilinearity}
Complete the proof of Proposition 7.3 by showing that $R(X,Y)(fZ) = fR(X,Y)Z$ for all smooth vector fields $X,Y,Z$ and smooth real-valued functions $f$.
\end{exercise}

\begin{exercise}
\label{prob:flat-connection}
\dcref{ch7:7-12}
Let $M$ be a smooth manifold and let $\nabla$ be any connection on $TM$. We can define the \textbf{curvature endomorphism} of $\nabla$ by the same formula as in the Riemannian case: $R(X,Y)Z = \nabla_X \nabla_Y Z - \nabla_Y \nabla_X Z - \nabla_{[X,Y]}Z$. Then $\nabla$ is said to be a \textbf{flat connection} if $R(X,Y)Z \equiv 0$. Prove that the following are equivalent:
\begin{enumerate}
\item[(a)] $\nabla$ is flat.
\item[(b)] For every point $p \in M$, there exists a parallel local frame defined on a neighborhood of $p$.
\item[(c)] For all $p,q \in M$, parallel transport along an admissible curve segment $\gamma$ from $p$ to $q$ depends only on the path-homotopy class of $\gamma$.
\item[(d)] Parallel transport around any sufficiently small closed curve is the identity; that is, for every $p \in M$, there exists a neighborhood $U$ of $p$ such that if $\gamma : [a,b] \to U$ is an admissible curve in $U$ starting and ending at $p$, then $P_{ab} : T_p M \to T_p M$ is the identity map.
\end{enumerate}
\end{exercise}

\begin{exercise}
\label{prob:kulkarni-nomizu-properties}
\dcref{ch7:7-9}
\uses{lem:kulkarni-nomizu-properties}
Prove parts (e) and (f) of Lemma 7.22 (properties of the Kulkarni--Nomizu product).
\end{exercise}

\begin{exercise}[Lichnerowicz's Theorem]
\label{prob:lichnerowicz-theorem}
\dcref{ch7:7-8}
\uses{prob:harmonic-functions,prob:eigenfunction-laplacian-inequality,prob:bochner-formula}
Suppose $(M, g)$ is a compact Riemannian $n$-manifold, and there is a positive constant $\kappa$ such that the Ricci tensor of $g$ satisfies $Rc(v, v) \geq \kappa|v|^2$ for all tangent vectors $v$. If $\lambda$ is any positive eigenvalue of $M$, then $\lambda \geq n\kappa/(n-1)$. [Hint: Use Probs. 2-23(c), 5-15, and 7-7.]
\end{exercise}

\begin{exercise}
\label{prob:lie-group-curvature}
\dcref{ch7:7-13}
\uses{prob:lie-group-bi-invariant-metric}
Let $G$ be a Lie group with a bi-invariant metric $g$. Show that the following formula holds whenever $X, Y, Z$ are left-invariant vector fields on $G$:
\[
R(X,Y)Z = \frac{1}{4}[Z,[X,Y]].
\]
\end{exercise}

\begin{exercise}
\label{prob:pseudo-riemannian-2d}
\dcref{ch7:7-15}
\uses{prop:canonical-form-vector-field}
Suppose $(M, g)$ is a 2-dimensional pseudo-Riemannian manifold of signature $(1, 1)$, and $p \in M$.
\begin{enumerate}
\item[(a)] Show that there is a smooth local frame $(E_1, E_2)$ in a neighborhood of $p$ such that $g(E_1, E_1) = g(E_2, E_2) = 0$.
\item[(b)] Show that there are smooth coordinates $(x, y)$ in a neighborhood of $p$ such that $(dx)^2 - (dy)^2 = fg$ for some smooth, positive real-valued function $f$. [Hint: Use Prop. A.45 to show that there exist coordinates $(t,u)$ in which $E_1 = \partial/\partial t$, and coordinates $(v,w)$ in which $E_2 = \partial/\partial v$, and set $x = u + w$, $y = u - w$.]
\item[(c)] Show that $(M, g)$ is locally conformally flat.
\end{enumerate}
\end{exercise}

\begin{exercise}
\label{prob:ricci-in-normal-coordinates}
\dcref{ch7:7-4}
Let $M$ be a Riemannian or pseudo-Riemannian manifold, and let $(x^i)$ be normal coordinates centered at $p \in M$. Show that the following holds at $p$:
\[
R_{ijkl} = \frac{1}{2}\left(\partial_j \partial_l g_{ik} + \partial_i \partial_k g_{jl} - \partial_i \partial_l g_{jk} - \partial_j \partial_k g_{il}\right).
\]
\end{exercise}
